% !TeX root = main.tex
% Shared sample: demonstrates numbered equations, figures, and a table inside
% an appendix. Included with \input from appendix.tex (single appendix) and
% appendixA.tex (lettered appendices); only one of those is compiled at a
% time, so each label is defined once. Delete this file and its \input lines
% when you write your own appendices.

\section{Equations and Figures in Appendices}

Equations inside the appendix are numbered with an appendix prefix:
\begin{align}
\label{eq:app1}
\nabla^2 \phi - \frac{1}{c^2}\frac{\partial^2 \phi}{\partial t^2} &= 0, \\
\label{eq:app2}
\oint_{\partial \Omega} \mathbf{F} \cdot d\mathbf{S} &= \iiint_{\Omega} (\nabla \cdot \mathbf{F})\, dV.
\end{align}
Cross-references like Equation~\eqref{eq:app1} and Equation~\eqref{eq:app2} resolve cleanly.

Figures placed within the appendix also receive appropriate appendix numbering, as shown in Figure~\ref{fig:app_cube1} and Figure~\ref{fig:app_cube2}.

\begin{figure}[htbp]
\centering
\includegraphics[width=2.5in]{msgcube}
\caption{First appendix demonstration figure.}
\label{fig:app_cube1}
\end{figure}

\begin{figure}[htbp]
\centering
\includegraphics[width=2.5in]{msgcube}
\caption{Second appendix demonstration figure.}
\label{fig:app_cube2}
\end{figure}

\clearpage
\section{Tables in an appendix}

Here is an example of a table in an appendix showing how you might typeset a table of websites you consulted in your thesis. 
Table~\ref{table:WEBSITES} lists common bioinformatics and signaling databases, formatted with automatic line wrapping for URLs using the \texttt{xurl} package.

\begin{table}[htbp]
\centering
\caption{Key Online Databases for Cell Signaling and Molecular Biology}
\label{table:WEBSITES}
\begin{tabularx}{\linewidth}{|c|p{0.32\linewidth}|X|}
\hline
\textbf{\#} & \textbf{Resource / URL} & \textbf{Description} \\
\hline
1 & \url{http://vlib.org/Science/Cell_Biology} & Virtual Library for Cell Biology with curated signal transduction pathways. \\
\hline
2 & \url{http://www.grt.kyushu-u.ac.jp/spad/} & Signaling Pathway Database containing curated pathway diagrams. \\
\hline
3 & \url{http://geo.nihs.go.jp/csndb/} & Cell Signaling Networks Database (CSNDB) \cite{CSND:1999}. \\
\hline
4 & \url{http://www.sdsc.edu/kinases/} & Protein Kinase Resource (PKR) for protein kinase structural data \cite{PKR:1997}. \\
\hline
5 & \url{http://www.expasy.ch/sprot/} & SWISS-PROT annotated protein sequence database \cite{BAIROCH:2000}. \\
\hline
6 & \url{http://www.expasy.ch/prosite/} & PROSITE database of protein families and functional domains \cite{HOFMANN:1999}. \\
\hline
7 & \url{http://www.rcsb.org/pdb/} & Protein Data Bank (PDB) for 3D macromolecular structures \cite{BERMAN:2000}. \\
\hline
\end{tabularx}
\end{table}
